\chapter{Εισαγωγή}
\label{chap:introduction}

\section{Αντικείμενο και κίνητρο}

Οι επιχειρήσεις ηλεκτρονικού εμπορίου συσσωρεύουν, εκ κατασκευής, πλήρες
ιστορικό των συναλλαγών τους. Κάθε παραγγελία, κάθε γραμμή προϊόντος και κάθε
επιστροφή καταγράφεται με ακρίβεια που σπανίως ήταν διαθέσιμη στο παραδοσιακό
λιανικό εμπόριο.

Η καταγραφή όμως δεν συνιστά γνώση. Στην πράξη, τα δεδομένα αυτά παραμένουν
συχνά αναξιοποίητα: χρησιμοποιούνται για λογιστικούς και λειτουργικούς σκοπούς,
ενώ οι αποφάσεις marketing και πωλήσεων λαμβάνονται με βάση την εμπειρία και
τη διαίσθηση. Το κενό ανάμεσα στα διαθέσιμα δεδομένα και στις αποφάσεις που θα
μπορούσαν να στηριχθούν σε αυτά αποτελεί το κίνητρο της παρούσας εργασίας.

Το αντικείμενο της εργασίας είναι η εξαγωγή αξιοποιήσιμων αναλύσεων και
προβλέψεων από δεδομένα αγοραπωλησιών, με τεχνικές εξόρυξης δεδομένων και
μηχανικής μάθησης, καθώς και η υλοποίηση ενός συστήματος που τις καθιστά
διαθέσιμες στην επιχείρηση.

\section{Στόχοι της εργασίας}

Οι αναλυτικοί στόχοι προέκυψαν από συζητήσεις με πραγματικούς πελάτες
ηλεκτρονικών καταστημάτων, και αφορούν ερωτήματα τα οποία οι ίδιοι δήλωσαν ότι
θα ήθελαν να απαντήσουν:

\begin{enumerate}
  \item Ποια αντικείμενα παρουσιάζουν τις μεγαλύτερες πωλήσεις.
  \item Ποιοι πελάτες πραγματοποιούν τις περισσότερες αγορές.
  \item Ποιες περιοχές παρουσιάζουν τις περισσότερες πωλήσεις.
  \item Ποια χρονική περίοδο συγκεντρώνονται οι πωλήσεις.
  \item Σε ποια προϊόντα καταγράφονται οι περισσότερες επιστροφές.
  \item Πώς κατανέμονται οι πωλήσεις ανά εύρος τιμής.
  \item Πώς κατανέμονται οι πωλήσεις ανά χαρακτηριστικό προϊόντος.
\end{enumerate}

Στα παραπάνω προστέθηκαν, ως επέκταση της αρχικής εκφώνησης, τα οικονομικά
μεγέθη που καθορίζουν αν οι απαντήσεις στα ερωτήματα αυτά αξίζει να
αξιοποιηθούν: η ακαθάριστη αξία πωλήσεων, η επαναληπτικότητα των πελατών, η
σύνθεση της ανάπτυξης και η κερδοφορία.

Πέραν των περιγραφικών ερωτημάτων, η εργασία εφαρμόζει τρεις οικογένειες
τεχνικών:

\begin{itemize}
  \item \textbf{Εξόρυξη συχνών υποσυνόλων}, για την ανάδειξη προϊόντων με
        ιδιαίτερη συσχέτιση, όπως αυτή αποτυπώνεται στα δεδομένα.
  \item \textbf{Συσταδοποίηση σε δύο επίπεδα}, με πρώτη ομαδοποίηση ως προς
        την περιοχή και τον χρόνο και δεύτερη, εντός κάθε συστάδας, ως προς
        επιπρόσθετες παραμέτρους συμπεριφοράς.
  \item \textbf{Τεχνικές πρόβλεψης}, για την προβολή των πωλήσεων στο μέλλον
        με μετρήσιμη ακρίβεια.
\end{itemize}

\section{Απαίτηση γενικότητας}

Δεύτερος, εξίσου ουσιώδης στόχος είναι η γενικότητα. Το σύστημα δεν σχεδιάστηκε
για ένα συγκεκριμένο κατάστημα αλλά ώστε να μπορεί να στραφεί σε οποιοδήποτε,
με ελάχιστη παραμετροποίηση.

Η απαίτηση αυτή διαπερνά τη σχεδίαση. Η ανάλυση δεν διαβάζει ποτέ το σχήμα
δεδομένων μιας συγκεκριμένης πηγής: κάθε πηγή προσαρμόζεται σε ένα
κανονικοποιημένο σχήμα, το οποίο δημοσιεύεται ως συμβόλαιο δεδομένων. Ένα
κατάστημα μπορεί να αποστείλει τις συναλλαγές του στη μορφή αυτή ή να
παραχωρήσει πρόσβαση ανάγνωσης στη βάση του, οπότε η προσάρτησή του αποτελεί
ζήτημα παραμετροποίησης και όχι συγγραφής κώδικα.

Για τον ίδιο λόγο αποφεύγονται τα απόλυτα κατώφλια όπου δεν δικαιολογούνται. Οι
δείκτες λογιστικής ανάπτυξης, για παράδειγμα, είναι εξ ορισμού ανεξάρτητοι από
το επιχειρηματικό μοντέλο, και η επιβολή ενός σταθερού σημείου αναφοράς θα
ακύρωνε ακριβώς αυτή την ιδιότητα.

\section{Ζήτημα ειλικρίνειας των αποτελεσμάτων}

Ένα σημείο διατρέχει ολόκληρη την εργασία και δηλώνεται εξαρχής: η διάκριση
μεταξύ όσων \emph{μετρώνται} και όσων \emph{υποτίθενται}.

Το σύνολο δεδομένων που χρησιμοποιείται καταγράφει τι πωλήθηκε και σε ποια
τιμή. Δεν καταγράφει κόστος πωληθέντων ούτε δαπάνες marketing. Συνεπώς
ερωτήματα κερδοφορίας δεν μπορούν να απαντηθούν ως μετρήσεις, όσο και αν η
απάντηση θα ήταν χρήσιμη.

Η εργασία δεν παρακάμπτει τον περιορισμό αυτόν με εύλογες εικασίες. Αντιθέτως,
τα μεγέθη κερδοφορίας υλοποιούνται ως ρητή ανάλυση ευαισθησίας πάνω σε
δηλωμένες παραδοχές, απομονωμένα από τα μετρήσιμα μεγέθη σε επίπεδο κώδικα,
διεπαφής και παρουσίασης. Η διάκριση αυτή αναπτύσσεται στην
ενότητα~\ref{sec:profitability-scenarios}.

\section{Δομή της εργασίας}

Το κεφάλαιο~\ref{chap:background} παρουσιάζει το θεωρητικό υπόβαθρο:
προεπεξεργασία δεδομένων, εξόρυξη συχνών υποσυνόλων, συσταδοποίηση, ανάλυση
RFM και τους οικονομικούς δείκτες του ηλεκτρονικού εμπορίου.

Το κεφάλαιο~\ref{chap:methodology} περιγράφει το σύνολο δεδομένων, την
αρχιτεκτονική του συστήματος, το κανονικοποιημένο σχήμα και τη διαδικασία
προεπεξεργασίας.

Τα κεφάλαια~\ref{chap:descriptive} έως~\ref{chap:forecasting-economics}
παρουσιάζουν τα αποτελέσματα: περιγραφική ανάλυση, συχνά υποσύνολα,
συσταδοποίηση, και τέλος πρόβλεψη μαζί με την οικονομική ανάλυση.

Το κεφάλαιο~\ref{chap:conclusions} συνοψίζει τα ευρήματα, τα αποτιμά ως προς
την πρακτική τους αξία, δηλώνει τους περιορισμούς και προτείνει επεκτάσεις.
